\chapter{Environment}
\section{Environmental Laws and policies}
\section{Climate Change, Carbon Credit, ETS}
\section{Environmental Convention and bodies}
\section{Pollution}
\section{Biodiversity}
\section{Protected Areas}
\section{Species}
\subsection{Rediscovery of Gynacantha khasiaca (Long-tailed Duskhawker)}

\begin{itemize}
    \item \textbf{Context \& Rediscovery:} 
    \begin{itemize}
        \item A rare dragonfly species, \textit{Gynacantha khasiaca} (commonly known as the long-tailed duskhawker), has resurfaced in the Changlang district of Arunachal Pradesh (\textbf{Namdapha National Park and Tiger Reserve}).
    \end{itemize}

    \item \textbf{About Dragonflies:} 
    \begin{enumerate}
        \item Possesses two large compound eyes, each comprised of thousands of tiny lenses and photoreceptor clusters.
        \item Features a near-$360^\circ$ field of vision and possesses the distinct behavioral trait of being able to hover perfectly still in mid-air.
        \item They serve as critical bioindicators and components of freshwater ecosystems, functioning simultaneously as apex invertebrate predators and vital prey within the aquatic food web.
        \item long-tailed duskhawker has also been historically documented across Assam, Meghalaya, West Bengal, Uttarakhand, and Maharashtra.
    \end{enumerate} 
\end{itemize}


\section{Miscellaneous}

\subsection{India's Green Transformation (2014--2026)}

\textbf{Three Pillars:} Ecological Capability $|$ National Capacity $|$ Global Credibility

\subsubsection*{Pillar 1: Ecological Capability}

\textbf{A. Forest Cover}
\begin{itemize}
  \item ISFR 2023: Forest + tree cover = 8.27 lakh sq. km (25.17\% of geographical area); forest cover alone = 7.15 lakh sq. km (21.76\%); carbon stock = 30.43 billion tonnes
  \item \textbf{Green India Mission (GIM)}: launched FY 2015--16; \rupee1019.26 cr released (till March 2026)
  \item \textbf{CAMPA} (Compensatory Afforestation Fund Act, 2016): 3.20 lakh ha afforested (FY21--25); GIS monitoring via HARIT-SANKALP
  \item \textbf{Nagar Van Yojana (NVY)}: 2020--27; target 1,000 Nagar Vans; 626 sanctioned; \rupee557.62 cr released
  \item \textbf{Aravalli Green Wall Initiative}: target 6.31 mn ha; 435 nurseries (393.24 lakh seedling capacity); 36,025 ha restored (2025)
  \item \textbf{Ek Ped Maa Ke Naam} (2024): 262.4 cr saplings planted (till Dec 2025); tracked via Meri LiFE portal
\end{itemize}

\textbf{B. River Rejuvenation -- Namami Gange Programme}
\begin{itemize}
  \item Launch: June 2014; initial outlay \rupee20,000 cr; Phase II (till March 2026): additional \rupee22,500 cr
  \item 524 projects (\rupee43,030 cr sanctioned); 355 completed; \rupee21,340 cr disbursed
  \item 218 STP projects (\rupee35,698 cr); total capacity 6,610 MLD; 138 STPs (3,977 MLD) operational
  \item Industrial BOD load: 26 TPD (2017) $\to$ 10.75 TPD (2024); effluent: 349 MLD $\to$ 265.56 MLD
  \item Afforestation: 33,024 ha; 7 Biodiversity Parks; 5 priority wetlands sanctioned
  \item 203 lakh IMC fingerlings released; 3,037 gharials (22 rivers); 6,327 Gangetic dolphins (28 rivers, 8,507 km surveyed)
  \item \textbf{Project Dolphin} (Aug 2020): Gangetic \& Indus Dolphins recognised as distinct species under WPA 1972; 200-km Dolphin Conservation Zone proposed (Chambal)
\end{itemize}

\textbf{C. Wetlands}
\begin{itemize}
  \item \textbf{NPCA} (National Plan for Conservation of Aquatic Ecosystems, 2013): 165 wetlands (incl. 42 Ramsar sites) sanctioned; cumulative release \rupee1,088.85 cr
  \item \textbf{Wetlands (Conservation \& Management) Rules, 2017}: prohibits encroachment, untreated effluent discharge, solid waste dumping
  \item Ramsar Sites: 26 (2014) $\to$ \textbf{99 (April 2026)}
\end{itemize}

\textbf{D. Mangroves}
\begin{itemize}
  \item \textbf{MISHTI Scheme} (Mangrove Initiative for Shoreline Habitats \& Tangible Incomes): initial outlay \rupee100 cr
  \item Mangrove cover: 4,628 sq. km (2013) $\to$ 4,992 sq. km (2023); net gain = 363 sq. km
\end{itemize}

\textbf{E. Coastal Ecosystems}
\begin{itemize}
  \item \textbf{National Coastal Mission}: extended 2025--31; allocation \rupee767 cr
  \item \textbf{CRZ Notification 2019}: science-based planning; Coastal Zone Management Plans (CZMPs) updated
  \item \textbf{Blue Flag Beaches}: 8 (2020) $\to$ 18 (2025--26); across 7 coastal states + 4 UTs
  \item \textbf{National Marine Turtle Action Plan (2021)}: 7,979 hatchlings released; 96.7\% hatching success; all 5 marine turtle species protected under WPA 1972
\end{itemize}

\textbf{F. Wildlife Conservation}
\begin{itemize}
  \item \textbf{Project Tiger} (est. 1973): Tiger reserves: 46 (2014) $\to$ 58 (2025); area $\approx$85,000 sq. km; tiger population: 2,226 (2014) $\to$ 3,682 (2022); India holds $>$70\% of world's wild tigers; 23 reserves with CA$|$TS accreditation (Dec 2023)
  \item \textbf{Project Cheetah} (Sep 2022): world's first intercontinental large carnivore translocation; 29 cheetahs brought (Namibia-8, S. Africa-12, Botswana-9); current population = 53
  \item \textbf{Project Lion} (Aug 2020): Asiatic lions: 523 (2015) $\to$ 891 (2025); range expanded $\approx$59\%
  \item \textbf{Leopard}: 12,852 (2018) $\to$ 13,874 (2022); NTCA + WII monitoring
  \item \textbf{Snow Leopard}: SECURE Himalaya (MoEFCC + UNDP + GEF, Oct 2017); SPAI 2019--23: 718 snow leopards (1,20,000 sq. km surveyed, $>$70\% habitat); SPAI 2.0 initiated
  \item \textbf{Project Elephant}: DNA-based All India Elephant Estimation (2021--25): 22,446 wild elephants; reserves: 26 (2014) $\to$ 33 (2025--26); corridors: 88 (2010) $\to$ 150 (2023); India holds $\approx$60\% of global wild Asian elephant population
  \item \textbf{One-Horned Rhinoceros}: National Conservation Strategy 2019; population: 1,500 (1980s) $\to$ 4,000+ (Sep 2024); $\approx$170\% increase
\end{itemize}

\subsubsection*{Pillar 2: National Capacity for Sustainable Transformation}

\textbf{A. Solid Waste Management}
\begin{itemize}
  \item Processing rate: 17\% (2014) $\to$ 77\% (2024); daily generation = 1,59,109 TPD; processed = 1,29,206 TPD
  \item Dumpsite remediation: 1,138 sites fully remediated (1,048 cities, 29 states); 877 lakh MT legacy waste cleared; 7,646 acres reclaimed
  \item \textbf{DRAP} (Dumpsite Remediation Accelerator Programme, Nov 2025): target zero dumpsites by Oct 2026
  \item \textbf{SWM Rules 2026} (supersedes 2016): 4-stream segregation (wet, dry, sanitary, special care); online tracking; \textbf{EBWGR} (Extended Bulk Waste Generator Responsibility); RDF substitution: 5\% $\to$ 15\% (over 6 years)
\end{itemize}

\textbf{B. Circular Economy -- EPR Framework (as of March 2026)}
\begin{itemize}
  \item 4,574 registered recyclers; total waste processed = 417.57 lakh MT; EPR certificates = 341.93 lakh MT
  \item Plastic: 196.97 lakh MT (2,986 recyclers); Battery: 69.37 lakh MT (520); Tyre: 122.29 lakh MT (579); E-Waste: 28.75 lakh MT (386); Used Oil: 0.19 lakh MT (103)
  \item Rules frameworks: Plastic WM Rules, E-Waste (Mgmt) Rules, Battery WM Rules, Waste Tyre EPR, Used Oil EPR, ELV framework, C\&D Waste Rules, EPR for non-ferrous scrap
\end{itemize}

\textbf{C. Green Skills \& Environmental Education}
\begin{itemize}
  \item \textbf{EEAT Scheme}: NGC, NNCP, CBP; 21 States + 1 UT; 1 lakh+ eco-clubs; 5.5 lakh students
  \item \textbf{EEARSD Scheme}: 1.34 cr students; 1.14 lakh eco-clubs (23 States); 4 lakh+ students in Eco-Creativity \& Innovation Hackathon
\end{itemize}

\textbf{D. Wildlife Science \& Technology}
\begin{itemize}
  \item \textbf{Next Generation DNA Sequencing facility} inaugurated at WII, Dehradun (Dec 2024): wildlife genetics, forensics, population studies
  \item \textbf{HAWK} (AI-based): real-time anti-poaching surveillance; \textbf{Gajah Suchana}: mobile platform for elephant data \& forensics
  \item WCCB: 166 joint operations in NE Region (2019--23); 375 arrests; cooperation with INTERPOL, SAWEN, CITES
  \item BSI \& ZSI: digitisation of herbarium \& faunal collections via dedicated databases
\end{itemize}

\textbf{E. Disaster Resilience}
\begin{itemize}
  \item \textbf{NDMP} (2016, revised 2019); 38 hazard-specific guidelines; Dynamic Composite Risk Atlas; database of 28,000 Himalayan glacial lakes
  \item \textbf{CAP-based Integrated Alert System}: geo-targeted multi-platform warnings; 4,500+ crore alert messages disseminated
  \item Apps: Damini, Mausam, Meghdoot; NDRF strengthened; 330+ institutions in Indian Universities \& Institutions Network for DRR
\end{itemize}

\textbf{F. Green Credit Programme (GCP)}
\begin{itemize}
  \item Launched under \textbf{Green Credit Rules, 2023}; aligned with Mission LiFE
  \item As of March 2026: 4,391 ha degraded forest land identified across 12 states for eco-restoration
\end{itemize}

\subsubsection*{Pillar 3: Global Credibility \& Climate Diplomacy}

\textbf{NDC Achievements}
\begin{itemize}
  \item Emissions intensity reduction: 33--35\% target met \textbf{11 years early}; actual reduction $>$36\% (from 2005 levels)
  \item 40\% non-fossil electricity capacity target met \textbf{9 years early}; non-fossil share = 52.57\% (Feb 2026)
  \item Additional carbon sink created: 2.29 billion tonnes CO$_2$ equivalent
  \item NDC 1.0 (2015) $\to$ NDC 2.0 (2022) $\to$ NDC 3.0 (2026): progressively enhanced 2030/2035 targets
\end{itemize}

\textbf{Key International Initiatives}
\begin{itemize}
  \item \textbf{ISA} (International Solar Alliance, 2015): launched India--France at COP21; 112 member countries
  \item \textbf{OSOWOG} (One Sun One World One Grid, 2018): India--UK Green Grids Initiative launched at COP26 (2021)
  \item \textbf{CDRI} (Coalition for Disaster Resilient Infrastructure, 2019): launched at UNGA Climate Summit; HQ New Delhi; status of international organisation (2022)
  \item \textbf{COP14-UNCCD} (New Delhi, 2019): Delhi Declaration; Land Degradation Neutrality by 2030; convergence of three Rio Conventions
  \item \textbf{CMS COP-13} (Gandhinagar, 2020): India as COP President; first CMS COP in India; norm-setter for migratory species, flyways, transboundary habitats
  \item \textbf{Mission LiFE} (2022): launched with UN Secretary-General; reflected in COP27 Sharm El-Sheikh Implementation Plan; mindful utilisation vs. mindless consumption
  \item \textbf{G20 New Delhi Leaders' Declaration -- Green Development Pact} (9 Sep 2023): India as consensus builder; Paris Agreement reaffirmation; climate finance, circular economy, biodiversity
  \item \textbf{IBCA} (International Big Cat Alliance, 2023): 7 big cat species; treaty-based intergovernmental org (23 Jan 2025); 26 member countries (Asia, Africa, Latin America)
  \item \textbf{WSDS 2024}: clean energy transition, climate justice, nature-based solutions
\end{itemize}






